\SourceLeaf{9}
\Article{2}
\par
Wir betrachten ganz allgemein einen endlichen Körper von beliebiger Gestalt, der vom übrigen unendlichen Raum durch eine zusammenhängende Fläche oder durch mehrere zusammenhängende, voneinander getrennte Flächen abgegrenzt ist (wenn der Körper einen oder mehrere Hohlräume einschließt). Die Gesamtheit dieser Flächen nennen wir einfach die Oberfläche des Körpers. Man denke sich diese Oberfläche in unendlich viele Elemente $\mathrm{d}s$ zerlegt. Es sei $P$ ein Punkt des Elements $\mathrm{d}s$, dessen Koordinaten, bezogen auf drei paarweise senkrechte Ebenen, mit $x,y,z$ bezeichnet werden. Es seien $PX,PY,PZ$ Geraden, die zu den jeweiligen Koordinatenachsen parallel und nach den Seiten gerichtet sind, zu denen die Koordinaten der Annahme nach positive Zuwächse erfahren. Ferner sei $PQ$ die nach außen gerichtete Normale der Oberfläche. Es sei $M$ der beliebig gelegene angezogene Punkt, seine Koordinaten seien $a,b,c$, und der Abstand $PM$ (stets positiv zu nehmen) sei $=r$. Die Winkel, die die Gerade $PM$ mit $PX,PY,PZ$ bildet, bezeichnen wir mit $MX,MY,MZ$ und die Winkel zwischen $PQ$ und $PX,PY,PZ,PM$ mit $QX,QY,QZ,QM$. All dies bezieht sich auf Punkte der Oberfläche im Allgemeinen. Wenn es sich um mehrere bestimmte Punkte der Oberfläche handelt, verwenden wir dieselben Zeichen mit unterschiedlichen Strichen.

\Article{3}
\par
Man denke sich eine Ebene senkrecht zur $x$-Koordinatenachse, und zwar so gelegen, dass, wenn man ihre Gleichung als $x=\alpha$ schreibt, $\alpha$ kleiner ist als der kleinste Wert der Koordinate $x$ auf der Oberfläche des Körpers. Die Projektion des Körpers auf diese Ebene ergibt dort eine endliche Figur, die wir uns in unendlich viele Elemente $\mathrm{d}\Sigma$ zerlegt denken. Im Element $\mathrm{d}\Sigma$ errichte man in einem Punkt $\Pi$ eine Senkrechte (also parallel zur $x$-Koordinatenachse), die den Körper in den Punkten $P',P'',P'''$ usw. schneidet. Die Zahl dieser Punkte ist offenbar gerade. Man errichte auch in jedem Randpunkt des Elements $\mathrm{d}\Sigma$ eine Senkrechte zur Ebene. Diese Senkrechten bilden eine Zylinderfläche im weiteren Sinne und schneiden aus der Oberfläche des Körpers die Elemente $\mathrm{d}s',\mathrm{d}s'',\mathrm{d}s'''$ usw. aus. Das Element $\mathrm{d}\Sigma$ ist die Projektion jedes einzelnen der Elemente $\mathrm{d}s',\mathrm{d}s'',\mathrm{d}s'''$ usw. Daraus folgt $\mathrm{d}\Sigma=\pm\mathrm{d}s'.\cos QX'=\pm\mathrm{d}s''.\cos QX''=\pm\mathrm{d}s'''.\cos QX'''$ usw., wobei das obere beziehungsweise untere Vorzeichen gilt, je nachdem, ob der Kosinus zu einem spitzen oder einem stumpfen Winkel gehört. Da aber die Senkrechte offenbar in $P'$ in den Körper eintritt, in $P''$ aus ihm austritt, in $P'''$ erneut eintritt und so weiter, sieht man leicht, dass $QX'$ stumpf, $QX''$ spitz, $QX'''$ stumpf ist und so weiter. Somit gilt
